\chapter*{Conclusion}

% rappel : problématique : 
%Dès lors, dans quelle mesure la variabilité du modèle de données FRBR adapté aux archives audiovisuelles conditionne-t-elle la mise en place d'une segmentation, et d'une indexation automatiques pour ce type d'archives ? Quels types d'algorithmes, fondés sur l'intelligence artificielle, peuvent être proposés au sein d'un outil de gestion des collections, à destination d'institutions patrimoniales ?


Le point de départ de cette étude est un constat d'inexploitabilité des fonds audiovisuels en raison de leur nature temporelle, qui en contraint la projection et nuit à l'efficacité de leur indexation documentaire. Si l'outil de gestion des collections permet de structurer et de faciliter la description via un modèle de données adaptatif, il ne résout pas encore le caractère chronophage des tâches d'indexation dédiées aux archives audiovisuelles. Le développement de l'intelligence artificielle offre des solutions intéressantes à ces difficultés, puisqu'elle permet d'automatiser l'indexation, qui suppose, dans le cas des archives audiovisuelles, une segmentation préalable. Ainsi nous avons cherché à répondre à la question suivante : dans quelle mesure la variabilité institutionnelle et sémantique du modèle FRBR adapté aux archives audiovisuelles limite-t-elle la mise en place d'une indexation automatique, et comment une intelligence hybride entre humain et machine permet-elle de dépasser ce verrou ?

lien indexation : FRBR. si l'indexation permet de normaliser la description, le FRBR la structure en différentes couches pour constituer , afin que l'utilisateur puisse se repérer au sein du logiciel

En amont d'une proposition d'un processus automatique, nous avons évoqué le modèle de données FRBR dans lequel s'inscrit l'indexation des archives audiovisuelles, conçu à l'origine pour les bibliothèques, et adapté aux images animées par la \gls{fiaf}. À travers l'analyse de la structure de données d'un corpus spécifique au sein de l'espace de travail S-Movies, développé par Skinsoft, nous avons évoqué la variabilité du FRBR et les complexités de navigation qu'il engendre. Dépendant de la nature du corpus décrit, le modèle de données, paramétré et adapté différemment pour chaque institution, conditionne ainsi l'implémentation de fonctionnalités innovantes pour les archives audiovisuelles, dont le besoin identifié est une segmentation automatique, préalable à leur indexation. la variabilité du FRBR ne peut pas être réduite lors de sa prise en compte dans un processus technique, ce processus doit tenter de s'inscrire au sein des cette variabilité. 

Dans ce cadre, nous avons proposé un processus de traitement automatique à partir de nos observations au sein du logiciel de gestion des collections développé par Skinsoft. Processus possible, à condition qu'il place la validation de l'utilisateur au centre et que le modèle de données soit repensé pour le reversement des données automatiquement produites. 
Face à l'absence de processus existants pour la segmentation des archives audiovisuelles en séquences, nous avons proposé d'adapter un outil de \gls{Vision par ordinateur} en libre accès, \gls{pyscene}, mis à l'épreuve à partir d'un corpus spécifique d'archives. Nous avons ensuite proposé une indexation automatique de ces séquences par mots-clés à l'aide du modèle \gls{yolo}. Si ce modèle est largement expoité, y compris dans le domaine patrimonial, la question de son entraînement reste décisive pour l'adapter à des corpus audiovisuels spécifiques. Le processus élaboré dans notre proposition résulte d'un choix d'une automatisation partielle qui place en son centre la validation d'un utilisateur, à la fois technique et métier. La place de l'utilisateur est notamment décisive pour assurer lea validation du reversement des données produites au sein du modèle de données existant, réadapté pour accueillir un tel processus. En effet, l'anticipation de la masse de données produites au cours de traitement automatiques nous a conduit à repenser le modèle de données, pour l'élargir à la notion de Fragment. L'étude du reversement des données automatiquement produite au sein du modèle de données existants, nous a conduit à proposer l'entité de "Fragment" proposée par l'IFLA, afin de rattacher les séquences segmentées à la hiérarchie existante, sans forcer le modèle dans des catégories inadaptées.

Dans cette optique, la proposition d'un tel processus s'inscrit finalement dans une structure fondée sur l'intelligence hybride réunissant l'homme et la machine, en cours de traitement, en amont et en aval, afin de valider la structure des données produites et leur qualité intrinsèque. Toutefois, elle reste conditionnée par les difficultés d'interface du système existants liées au modèle de données aux relations complexes, qu'il s'agit d'améliorer afin d'instaurer un tel processus, centré sur la validation humaine. L'amélioration de l'interface permettrait d'intégrer au mieux l'utilisateur en amont du processus, pour produire des données d'entrainment, et en aval, afin de valider la qualité et la structure des données produites. La notion d'intelligence hybride répond notamment aux enjeux d'entraînement du modèle \gls{yolo}, qui pourrait mettre à profit les données audiovisuelles constituées par les clients de Skinsoft. Cette notion répond également aux difficultés de généralisation du processus. En proposant un processus paramétrable, cela permet d'intégrer l'utilisateur métier dès le début et de l'intégrer aux besoins nécessaires quant au modèle de données et au processus à effectuer. 
les perspectives ouvertes par l'IA sont nombreuses et placent l'éditeur logiciel au sein d'un rôle pivot, facilitateur de fonctionnalités innovantes pour aider la découvrabilité des fonds audiovisuels, tout en adaptant aux besoins patrimoniaux documentaire d'interopérabilité des données : adapter ces fonctionnalités au FRBR. 
le rôle d'interface entre des outils technologiques génériques et des institutions patrimoniales aux besoins spécifiques. Ni l'automatisation totale promise par les grandes plateformes commerciales, ni le maintien du tout-manuel des pratiques traditionnelles — mais une voie intermédiaire, adaptative, centrée sur l'utilisateur et respectueuse de la complexité documentaire des collections qu'elle sert.



%La troisième partie a évalué les limites structurelles de cette démarche. Trois limites se sont révélées particulièrement robustes. La sous-représentation des archives audiovisuelles patrimoniales historiques dans les corpus d'entraînement disponibles est la plus fondamentale : les modèles sont entraînés sur des vidéos contemporaines en couleur, et leurs performances se dégradent significativement sur des archives muettes en noir et blanc. Le fine-tuning et le transfer learning offrent des pistes, mais supposent des ressources humaines et computationnelles que la plupart des institutions ne possèdent pas. La variabilité institutionnelle du FRBR constitue une limite complémentaire : même si un modèle générique peut produire des métadonnées pertinentes, leur insertion correcte dans la hiérarchie documentaire reste une opération qui exige une connaissance des règles métier propres à chaque institution. Enfin, la question éthique de la transparence s'impose : la data provenance — la traçabilité de chaque métadonnée automatique jusqu'à sa source algorithmique — est une condition de confiance sans laquelle l'archiviste ne peut ni valider intelligemment les propositions du modèle, ni assumer la responsabilité des données intégrées dans la base.

%et la manière dont il conditionne l'implémentation de nouvelles fonctionnalités. Dans le cadre de ce mémoire de recherche, nous avons mesuré la complexité de la mise en place d'une segmentation puis d'une indexation automatisées pour les collections audiovisuelles, alors même que le modèle de données contraint la navigation et la compréhension documentaire pour les utilisateurs. 
%L'analyse de l'adaptation du modèle \gls{FRBR} aux archives audiovisuelles a notamment permis de pointer sa verticalité, 
%malgré une adaptation par la \gls{fiaf}, pose des difficultés pour structurer une interface de gestion pour des corpus audiovisuels spécifiques. Concu pour organiser la description et facilite la navigation, il pose en réalité des problèmes d'interface en permettant de multiplier les liens, supposant une interface innovante qui permet de aviguer et de comprendre les liens et la place des entités au niveau hiérarchique. Difficile d'adapter un modèle conçu par les bibliothèques aux archives audiovisuelles qui sont hétérogènes et reproductibles : multiplicité des liens (segments, extraits, bobines, numériques...). nécessite une interface 
%Ainsi, l'utilisation du FRBR et la navigation entre ses différents niveaux dépend de la nature du corpus audiovisuel et est, de fait, hétérogène.  
%C'est dans ce contexte que nous avons placé l'étude d'un projet de mise en place d'une segmentation automatique des archives audiovisuelles, puis de leur indexation par mots-clés elle aussi automatisées, projet pensé par l'éditeur Skinsoft pour une demande institutionnelle. Ce projet, toujours à l'état de projet et de réflexions, n'est pas encore développé. Notre étude est une proposition et une mise à l'épreuve d'outils, en lien avec les enjeux de modélisation et d'interface, à travers l'exemple d'un corpus spécifique afin d'en déceler les difficultésn les enjeux propres aux archives audiovisuelles. 
%Étant donné l'absence de solutions de segmentation automatiques en séquences, nous avons proposé la mise à l'épreuve d'un outil de vision par ordinateur classique conçu pour la segmentation d'un flux vidéos en ses différents plans. Puis les segments obtenus sont ensuite indexés par des mots clés par un outil de vision par ordinateur fondé sur des réseaux de neurones \gls{cnn}. Nous avons analysé les modalités d'netraînement d'un tel modèle, sur des images d'archives. Pour que ce processus automatique reste une assistance à l'usager, une aide (enjeux de transparence et de production qualité des données), nous avons imaginé une manière d'intégrer l'utilisateur au processus à chaque étape, donc un processus de validation. afin de valider les données produites puis de les reverser adéquatement au sein du modèle de données. Nous avons prévu une modification du modèle de donénes, face à la massification anticipée des données créées automatiquement. 
%
%Dès lors, nous en avons conclu que la modélisation du FRBR selon le contexte institutionnel constitue un frein à la mise en place de l'indexation automatique, puisqu'elle repose sur des règles métier d'usage selon le corpus indexé. La refonte de la navigation entre les niveau du FRBR constitue alors un tournant essentiel et nécessaire à la mise en place d'une telle automatisation, permettant aux modèles d'apprentissage et à l'utilisateur de mieux comprendre les liens entre les différents niveaux d'indexation. 
%
%Partie 3 
%
%Mise en place d'une intelligence hybride et contrôle qualité + nécessité d'evolution de l'espace de travail en accord avec la nécessité de garder une approche d'ergonomie d'interface comme les applications web commerciales. Toutefois, notre étude est une proposition, c'est pourquoi nous avons proposé une mise en place expérimentale à travers un exemple. complexités de généralisation du processus à l'ensemble de l'institutionpose également la question des besoins de telles foncionnalités. mise en place de fonctionnalités modulables, paramétrables, s'adaptant ainsi au FRBR

%Finalement, il s'agit de centrer notre approche sur l'expérience utilisateur en raison des difficultés d'entraîner un modèle d'intelligence artificielle adapté aux corpus audiovisuels et au modèle de données associé. Nous proposons alors la mise en place d'une aide à l'indexation, d'une indexation augmentée, qui suggèrerait à l'utilisateur des choix en les explicitant, sans pour autant automatiser l'ensemble du processus.  Par ses choix de validation, l'utilisateur entraînerait lui-même le modèle à la fois sur les corpus utilisés et à la fois sur l'utilisation du modèle FRBR. Une telle interface d'échange pourrait permettre de rendre l'indexation plus efficace et d'améliorer la compréhension des liens entre les niveaux d'indexation du FRBR. 

